%% sections/validity.tex
%% \owner{ANCHOR subagent}  GOAL: fill REAL construct-validity numbers and an honest
%% reading of them. Currently the committed bench/out/anchor.json has radon/lizard/
%% cognitive as SKIP (only Buse--Weimer ran, aggregate Spearman ~0.39). The subagent must:
%%   1. Run anchor.py in the quarantined venv so radon/lizard/cognitive are REAL:
%%        python3 -m venv /tmp/anchorvenv
%;        /tmp/anchorvenv/bin/pip install radon lizard cognitive_complexity
%%        /tmp/anchorvenv/bin/python bench/anchor.py    # -> bench/out/anchor.json
%%   2. Extract the headline numbers (Spearman of each metric vs the incidental knob rank)
%%      and fill Table~\ref{tab:validity} below.
%%   3. Write the honest reading: complexity metrics (radon/lizard CC) move strongly with
%%      the knob BUT that is partly BY CONSTRUCTION (the knob is built to inflate them);
%%      the less-tautological readability signal (Buse--Weimer features) is WEAKER
%%      (aggregate ~0.39) and a human-calibrated readability study is future work. Do NOT
%%      overclaim. This honesty is itself a selling point for a data paper.

\section{Construct Validity: do the labels mean anything?}
\label{sec:validity}

The by-construction quality order is only useful if it tracks quantities the community
already trusts. We test whether the incidental knob moves established static metrics
monotonically, over the public \texttt{dev} set ($50$ base samples expanded to $300$
sources across the six knob levels). For each base sample we hold semantics fixed and
compute Spearman's $\rho$ between the knob rank
(\texttt{clean}${}\prec{}$\texttt{minimal}${}\prec{}$\texttt{light}${}\prec{}$%
\texttt{standard}${}\prec{}$\texttt{heavy}${}\prec{}$\texttt{max}) and each metric,
reporting the mean over samples (Table~\ref{tab:validity}). The anchor tools
(\texttt{radon}~6.0.1, \texttt{lizard}~1.23.0, \texttt{cognitive\_complexity}~1.3.0)
are kept off the zero-dependency core and run only inside a quarantined virtualenv; a
missing tool records an honest \texttt{SKIP} rather than a fabricated number.

\begin{table}[t]
\centering
\caption{Spearman correlation between the incidental ``messiness'' knob rank and
established static metrics, averaged over the $50$ \texttt{dev} base samples (higher
$|\rho|$ = the label order is better reflected). Complexity metrics move strongly with
the knob, but partly \emph{by construction}; the less tautological readability signal is
markedly weaker.}
\label{tab:validity}
\begin{tabular}{@{}llr@{}}
\toprule
Metric (tool) & Direction & Spearman $\rho$ vs.\ knob \\
\midrule
Cyclomatic complexity (radon)        & + & $+0.75$ \\
Cyclomatic complexity (lizard)       & + & $+0.75$ \\
Cognitive complexity                 & + & $+0.80$ \\
Maintainability index (radon)        & $-$ & $-0.79$ \\
Buse--Weimer surface density (agg.)  & + & $+0.39$ \\
\bottomrule
\end{tabular}
\end{table}

\paragraph{Complexity moves strongly, but read this with care.}
The complexity anchors agree and move sharply with the knob: cyclomatic complexity
(identical at $\rho=+0.75$ under both \texttt{radon} and \texttt{lizard}), cognitive
complexity ($\rho=+0.80$), and the maintainability index in the expected
\emph{opposite} direction ($\rho=-0.79$). Two independent tool families converging on
the same ordering, and our own stdlib metric lane cross-checking against them at
$\rho>0.97$, is reassuring evidence that the labels are not arbitrary. We deliberately
under-claim here, though, because this agreement is partly \emph{tautological}: the
incidental knob is engineered to inject redundant control flow, dead branches, and
nested scaffolding, which is precisely what cyclomatic and cognitive complexity count.
A metric rising when we add the very constructs it was designed to detect confirms that
the transformation does what it says, but it is weaker evidence that the order matches
\emph{human} judgments of quality.

\paragraph{The readability signal is real but weaker; human calibration is future work.}
The more informative test is therefore the one anchor that is \emph{not} a complexity
counter. Our re-implementation of the Buse--Weimer~\citep{buse2010readability} surface
readability feature set (computed with the Python stdlib over the same sources)
correlates with the knob at only $\rho=+0.39$ in the documented density aggregate ---
roughly half the strength of the complexity metrics. We read this honestly: the knob
makes code denser and harder to read, but far less unambiguously than it inflates
control-flow complexity, and even this number is a surface-density proxy rather than the
original trained readability classifier (whose coefficients are not available offline,
so that anchor remains a recorded \texttt{SKIP}). The benchmark's construct validity
thus rests on convergent static evidence, not on a claim of calibrated human
readability. A fresh human-rated readability study aligned to this exact stimulus set is
the natural next step and is left as explicit future work; until then we treat the
$\rho=+0.39$ readability signal, not the $\rho=+0.80$ complexity signal, as the
conservative bound on how much the by-construction order reflects perceived quality.
